\documentclass[11pt,a4paper]{article}
\usepackage[utf8]{inputenc}
\usepackage[T1]{fontenc}
\usepackage{amsmath,amssymb,amsfonts}
\usepackage{booktabs}
\usepackage{multirow}
\usepackage{graphicx}
\usepackage{xcolor}
\usepackage{hyperref}
\usepackage[margin=2.5cm]{geometry}
\usepackage{natbib}
\usepackage{tabularx}
\usepackage{caption}
\usepackage{subcaption}

\newcommand{\sica}{\textsc{Sica}}
\newcommand{\scmethod}{\textsc{SC}}
\newcommand{\czero}{$C_0$}

\title{The Capability Trap: Why Self-Extracted Constraints\\Cannot Improve LLM Majority Voting}

\author{Anonymous}
\date{}

\begin{document}
\maketitle

% ============================================================
% ABSTRACT
% ============================================================
\begin{abstract}
Recent work augments Large Language Model (LLM) reasoning with symbolic methods that extract first-order logic (FOL) constraints from reasoning traces and aggregate them via solvers such as MAX-SAT.
We ask a practical question: when does such symbolic aggregation improve over simple majority voting (Self-Consistency, SC)?
We instantiate the paradigm as \sica{} (Symbolic Invariant Consensus Aggregation), which extracts FOL constraints from sampled reasoning traces and selects answers via weighted MAX-SAT, and evaluate it across three reasoning domains (formal logic, commonsense, science) and five model families spanning instruction-tuned, distilled-reasoning, and full-reasoning architectures.
Across all 18 tested conditions (combined $N{>}1{,}300$), \sica{} provides no significant accuracy gain over SC, with per-condition $\Delta$ ranging from $-3.3$pp to $+2.0$pp.
We identify the structural cause: a \emph{capability trap} in which self-extracted constraints achieve precision $\approx 0.50$ (no better than chance) at distinguishing correct from incorrect answers, causing MAX-SAT to degenerate to weighted counting.
A constraint ablation confirms that MAX-SAT, count-only, and random-weight aggregation produce statistically identical results ($>$99.6\% agreement).
We propose $C_0$ (multi-candidate rate) as a zero-cost diagnostic: $C_0{<}5\%$ predicts method degeneracy, while high $C_0$ is necessary but insufficient for gains.
These findings provide practitioners with a principled early-stopping criterion before investing in symbolic aggregation overhead.
\end{abstract}

% ============================================================
% 1. INTRODUCTION
% ============================================================
\section{Introduction}
\label{sec:intro}

Self-Consistency (SC) \citep{wang2023selfconsistency} improves LLM reasoning by sampling multiple reasoning traces and selecting the majority answer.
A natural extension is to extract richer structure from these traces: formal constraints, logical invariants, or symbolic representations that a solver can aggregate more intelligently than simple voting \citep{pan2023logiclm,olausson2023linc}.
Several recent systems pursue this direction, using FOL translation, constraint satisfaction, or symbolic verification to rerank or select answers \citep{pan2023logiclm,olausson2023linc,dylos2026,aristotle2025,fover2025}.
The appeal is clear: if reasoning traces contain more information than just the final answer, a solver that exploits that information should outperform voting.

But when does this actually work?
Prior systems report gains on specific benchmarks, yet use trained verifiers, external oracles, or multi-paradigm solver selection.
A fundamental question remains open: can the \emph{same} LLM that generates reasoning traces also extract constraints of sufficient quality to improve its own voting?
We call this the \emph{self-extraction} setting, where a single model both reasons and formalizes.

We instantiate self-extraction as \sica{} (Symbolic Invariant Consensus Aggregation): sample $K$ reasoning traces, extract FOL constraints from each via the same LLM, and solve a weighted MAX-SAT instance to select the best answer.
We evaluate \sica{} across three reasoning domains (FOLIO, StrategyQA, ARC-Challenge, with GSM8K and ProofWriter as supplementary) and five model families (Qwen-14B, LLaMA-8B, Mistral-7B, DeepSeek-R1-Distill-7B, QwQ-32B), spanning instruction-tuned, distilled-reasoning, and full-reasoning architectures.

Our central finding is negative: \sica{} provides no significant improvement over SC in any tested condition.
We identify the structural cause as a \emph{capability trap}:
\begin{itemize}
    \item \sica{} requires output diversity (multiple candidate answers) to differ from voting. We quantify this via $C_0$, the fraction of problems with more than one candidate answer across $K$ traces.
    \item When $C_0$ is low ($<$5\%, as with full reasoning models), \sica{} structurally degenerates to SC because there is only one candidate.
    \item When $C_0$ is high (38--62\%, as with instruction-tuned models), self-extracted constraints achieve precision $\approx 0.50$ at discriminating correct from incorrect answers, making MAX-SAT equivalent to counting.
    \item No tested condition achieves both sufficient diversity AND discriminative constraints simultaneously.
\end{itemize}

Our contributions are:
\begin{enumerate}
    \item \textbf{Systematic negative result}: We show that self-extracted FOL constraint aggregation via MAX-SAT provides no significant accuracy gain over majority voting across 3 domains, 5 model families, and 18 experimental conditions (combined $N{>}1{,}300$).
    \item \textbf{Capability trap characterization}: We identify and quantify the structural failure mode: self-extracted constraints achieve chance-level discriminative precision ($\bar{p}=0.50$, $n=307$ multi-candidate problems), causing MAX-SAT to degenerate to counting (confirmed by three-way ablation: MAX-SAT $\equiv$ count-only $\equiv$ random-weight, $>$99.6\% agreement).
    \item \textbf{$C_0$ deployment diagnostic}: We propose the multi-candidate rate $C_0$ as a zero-cost early-stopping criterion. $C_0{<}5\%$ reliably predicts degeneracy; high $C_0$ is necessary but insufficient for gains (Pearson $r=0.05$, $p=0.84$ across 18 conditions).
\end{enumerate}

% ============================================================
% 2. RELATED WORK
% ============================================================
\section{Related Work}
\label{sec:related}

\paragraph{Self-Consistency and extensions.}
SC \citep{wang2023selfconsistency} established majority voting over sampled reasoning traces as a strong baseline for LLM reasoning.
Universal SC (USC) \citep{chen2023usc} extends this to free-form answers.
\citet{loo2025sc} demonstrate that SC scaling plateaus with increased samples, motivating richer aggregation.
Confidence-weighted SC \citep{taubenfeld2025weighted} assigns trace-level weights but remains within the voting paradigm.
DiVeRSe \citep{li2023diverse} combines diverse prompting with verification.

\paragraph{Neuro-symbolic reasoning.}
Logic-LM \citep{pan2023logiclm} and LINC \citep{olausson2023linc} translate natural language to FOL for symbolic solving, documenting that LLM-to-FOL translation is error-prone.
Recent systems achieve gains using trained verifiers or multi-paradigm approaches: DyLoS \citep{dylos2026} uses dynamic logic-solver selection, Aristotle \citep{aristotle2025} employs logic-complete decomposition, and FoVer \citep{fover2025} trains a dedicated verifier.
Our work differs by studying the specific failure mode of \emph{same-model self-extraction} pipelines rather than these richer architectures.

\paragraph{Verification and reranking.}
Outcome Reward Models (ORM) \citep{cobbe2021gsm8k} and Process Reward Models (PRM) \citep{lightman2023prm} use trained verifiers to rerank solutions.
Math-Shepherd \citep{wang2024mathshepherd} extends process supervision.
These approaches train separate verifiers, sidestepping the self-extraction bottleneck we characterize.

\paragraph{Test-time compute scaling.}
\citet{snell2025scaling} analyze optimal test-time compute allocation.
DeepSeek-R1 \citep{deepseekr1} and OpenAI o1 \citep{openai2024o1} demonstrate that extended reasoning chains can substitute for external verification.
Our finding that full reasoning models (QwQ-32B) produce near-zero diversity ($C_0{<}5\%$) is consistent with this paradigm: models that ``think longer'' converge to consistent answers, eliminating the operating regime for symbolic aggregation.

\paragraph{Constraint satisfaction in NLP.}
SATBench \citep{satbench2025} evaluates LLM performance on satisfiability problems.
Recent work explores constraint propagation with LLMs \citep{llmcp2025}.
Our MAX-SAT formulation draws on this line but applies it to answer aggregation rather than direct problem solving.

% ============================================================
% 3. METHOD
% ============================================================
\section{Method: SICA}
\label{sec:method}

\sica{} (Symbolic Invariant Consensus Aggregation) aggregates LLM reasoning traces via symbolic constraint satisfaction.
The pipeline has four stages.

\paragraph{Stage 1: Trace generation.}
Given a problem $q$, sample $K$ reasoning traces $\{t_1, \ldots, t_K\}$ from an LLM with temperature $\tau > 0$.
Extract the final answer $a_i$ from each trace $t_i$.
Let $\mathcal{A} = \{a_1, \ldots, a_m\}$ be the set of distinct candidate answers.

\paragraph{Stage 2: Constraint extraction.}
For each trace $t_i$, prompt the same LLM to extract a set of FOL constraints $\mathcal{C}_i = \{c_{i,1}, \ldots, c_{i,n_i}\}$ that formalize key reasoning steps.
Each constraint $c_{i,j}$ is a FOL formula intended to capture an invariant of the reasoning process.

\paragraph{Stage 3: Constraint evaluation.}
For each candidate answer $a \in \mathcal{A}$ and each constraint $c_{i,j}$, evaluate whether $a$ satisfies $c_{i,j}$.
Define $\text{sat}(a, c_{i,j}) \in \{0, 1\}$.
This produces a satisfaction matrix of size $|\mathcal{A}| \times |\mathcal{C}|$ where $\mathcal{C} = \bigcup_i \mathcal{C}_i$.

\paragraph{Stage 4: MAX-SAT aggregation.}
Formulate a weighted MAX-SAT instance where each constraint-answer pair contributes a weighted clause.
The answer maximizing total satisfied constraint weight is selected:
\begin{equation}
    a^* = \arg\max_{a \in \mathcal{A}} \sum_{c \in \mathcal{C}} w_c \cdot \text{sat}(a, c)
\end{equation}
where $w_c$ is the weight assigned to constraint $c$ (default: uniform).
In case of ties, \sica{} selects the answer supported by the trace with the highest individual constraint satisfaction count.

\paragraph{Baseline: Self-Consistency (SC).}
SC selects $a^* = \arg\max_{a \in \mathcal{A}} |\{i : a_i = a\}|$, using only final answers.

\paragraph{$C_0$ diagnostic.}
We define $C_0$ as the fraction of problems where $|\mathcal{A}| > 1$ (i.e., traces produce more than one distinct answer).
When $C_0 = 0$, all traces agree and \sica{} trivially equals SC.
$C_0$ is computable from Stage 1 alone, before any constraint extraction, making it a zero-cost diagnostic.

% ============================================================
% 4. EXPERIMENTS
% ============================================================
\section{Experiments}
\label{sec:experiments}

\subsection{Setup}
\label{sec:setup}

\paragraph{Domains.}
We evaluate on three primary domains: \textbf{FOLIO} (formal deductive logic, 204 problems), \textbf{StrategyQA} (commonsense reasoning, 200 problems for powered tests), and \textbf{ARC-Challenge} (science reasoning, 60 problems).
Supplementary domains include \textbf{GSM8K} (arithmetic, 200 problems) and \textbf{ProofWriter} (included in QwQ-32B evaluation).

\paragraph{Models.}
We test five model families:
(1) \textbf{Qwen2.5-14B-Instruct} as the primary model ($n=804$ powered evaluation across domains);
(2) \textbf{LLaMA-3.1-8B-Instruct} on StrategyQA ($n=200$) and FOLIO ($n=204$);
(3) \textbf{DeepSeek-R1-Distill-Qwen-7B} ($n=60$);
(4) \textbf{QwQ-32B} as a full reasoning model ($n=46$);
(5) Within-family scaling with \textbf{Qwen2.5-7B} and \textbf{Qwen2.5-72B-AWQ}.

All experiments use $K=12$ sampled traces unless otherwise noted.

\paragraph{Metrics.}
We report accuracy for both \sica{} and SC, their difference $\Delta$ (in percentage points), McNemar's test $p$-value for significance, and 95\% bootstrap confidence intervals for powered tests.
We also report $C_0$ (multi-candidate rate), constraint extraction rate, and constraint precision (fraction of constraints that correctly distinguish the gold answer from incorrect candidates).

% ============================================================
% 4.2 MAIN RESULTS
% ============================================================
\subsection{Main Results}
\label{sec:main_results}

Table~\ref{tab:main} summarizes results across all model-domain combinations.
Across every tested condition, \sica{} provides no statistically significant improvement over SC.

\begin{table*}[t]
\centering
\caption{Main results: \sica{} vs.\ SC across models and domains. $\Delta$ = \sica{} $-$ SC in percentage points. $C_0$ = multi-candidate rate. All $p$-values from McNemar's test. No condition achieves significance at $\alpha=0.05$.}
\label{tab:main}
\small
\begin{tabular}{llrrrrrl}
\toprule
\textbf{Model} & \textbf{Domain} & \textbf{$n$} & \textbf{\sica{} Acc} & \textbf{SC Acc} & \textbf{$\Delta$pp} & \textbf{$C_0$\%} & \textbf{$p$} \\
\midrule
Qwen2.5-14B & FOLIO+SQA+ARC & 804 & 76.5 & 76.2 & +0.25 & 38 & -- \\
Qwen2.5-14B & StrategyQA & 60 & 56.7 & 53.3 & +3.33 & 10 & 0.50 \\
Qwen2.5-14B & ARC-Challenge & 60 & 95.0 & 96.7 & $-$1.67 & 5 & 1.00 \\
Qwen2.5-14B & GSM8K & 30 & 86.7 & 86.7 & 0.0 & 13 & -- \\
Qwen2.5-14B & GSM8K (expanded) & 200 & 89.0 & 89.0 & 0.0 & 12 & 1.00 \\
\midrule
LLaMA-3.1-8B & StrategyQA & 200 & -- & -- & +2.0 & 62 & 0.424 \\
LLaMA-3.1-8B & FOLIO & 204 & -- & -- & $-$2.45 & 62 & 0.267 \\
\midrule
DeepSeek-R1-7B & D012 (pilot) & 60 & 68.3 & 71.7 & $-$3.33 & 58 & 0.48 \\
\midrule
Qwen2.5-7B & D012 (pilot) & 60 & 58.3 & 58.3 & 0.0 & 70 & 1.00 \\
Qwen2.5-72B-AWQ & D012 (pilot) & 60 & 80.0 & 80.0 & 0.0 & 22 & 1.00 \\
\midrule
QwQ-32B & All domains & 46 & -- & -- & 0.0 & 2.6 & -- \\
\bottomrule
\end{tabular}
\end{table*}

Three patterns emerge from Table~\ref{tab:main}:

\textbf{Instruction-scale models show high diversity but zero gain.}
Qwen-14B on its primary evaluation ($n=804$) achieves $C_0=38\%$, meaning over a third of problems produce multiple candidate answers.
Yet $\Delta=+0.25$pp, well within noise.
LLaMA-8B shows the highest diversity ($C_0=62\%$) with powered tests ($n\geq200$) yielding $\Delta=+2.0$pp on StrategyQA (CI: [$-1.5$, $+5.5$]) and $\Delta=-2.45$pp on FOLIO (CI: [$-5.88$, $+0.98$]), both non-significant and in opposite directions.
The meta-analytic $\Delta$ across LLaMA-8B conditions is $+0.3$pp.

\textbf{Full reasoning models produce near-zero diversity.}
QwQ-32B achieves an effective $C_0$ of 2.6\% across all domains (AIME, ProofWriter, FOLIO), with $\Delta=0.0$pp in every domain.
Extended reasoning chains produce consistent answers, structurally eliminating the operating regime for \sica{}.

\textbf{Within-family scaling confirms the trap from both ends.}
In the Qwen family, scaling from 7B ($C_0=70\%$, $\Delta=0.0$) through 14B ($C_0=38\%$, $\Delta=+0.25$pp) to 72B ($C_0=22\%$, $\Delta=0.0$) traces a capability-diversity tradeoff: weaker models produce more diversity but worse constraints, while stronger models produce better constraints but less diversity.
Neither end benefits from \sica{}.

% ============================================================
% 4.3 CONSTRAINT ABLATION
% ============================================================
\subsection{Constraint Ablation: MAX-SAT Adds Nothing Beyond Counting}
\label{sec:ablation}

If constraints carry discriminative information, the MAX-SAT solver should outperform simpler aggregation.
We test this by comparing three aggregation strategies on the Qwen-14B $n=804$ evaluation (Table~\ref{tab:ablation}).

\begin{table}[t]
\centering
\caption{Constraint ablation ($n=804$, Qwen-14B). Three aggregation methods produce statistically identical results, confirming that MAX-SAT adds nothing when constraints lack discriminative power.}
\label{tab:ablation}
\small
\begin{tabular}{lrr}
\toprule
\textbf{Aggregation} & \textbf{Accuracy} & \textbf{Agreement w/ MAX-SAT} \\
\midrule
MAX-SAT (full \sica{}) & 76.49 & -- \\
Count-only & 76.49 & 803/804 \\
Random-weight & 76.37 & 801/804 \\
SC (majority vote) & 76.24 & -- \\
\bottomrule
\end{tabular}
\end{table}

MAX-SAT and count-only agree on 803 of 804 problems (McNemar $p=1.0$).
Random-weight agrees on 801/804 (McNemar $p=0.82$).
The three-way equivalence is the mechanistic signature of the capability trap: when constraint precision $\approx 0.50$, each answer's satisfied constraint count is approximately proportional to its vote count.
The MAX-SAT solver, despite its computational machinery, reduces to counting, which reduces to voting.

% ============================================================
% 4.4 PER-PROBLEM CONSTRAINT ANALYSIS
% ============================================================
\subsection{Per-Problem Constraint Analysis}
\label{sec:constraint_analysis}

To understand \emph{why} constraints fail to discriminate, we analyze constraint precision at per-problem granularity across $n=307$ multi-candidate problems from the Qwen-14B evaluation.

\paragraph{Constraint precision is indistinguishable from random.}
Mean constraint precision is $0.50$ ($\sigma=0.12$), with 89\% of problems falling in the $[0.4, 0.6]$ range.
There is zero correlation between problem-level diversity and constraint quality (Pearson $r=0.03$, $p=0.58$).
Even among the 34 problems with precision $> 0.6$, \sica{} wins only 56\% of the time (barely above chance).

\paragraph{Constraints between positive and negative flips.}
Problems where \sica{} flips from incorrect to correct (positive flips) have mean constraint precision $0.52$; problems where \sica{} flips from correct to incorrect (negative flips) have precision $0.48$.
The difference is negligible, confirming that constraints do not systematically identify correct answers.

\paragraph{Qualitative analysis.}
We examine 8 representative problems in detail.
Of extracted constraints, 75\% are semantically meaningful FOL formulas, and 52\% correctly evaluate against their source answer.
However, only 12\% are \emph{discriminative}, producing different satisfaction values for different candidate answers.
The typical failure mode is tautological constraints: a formula like $\forall x (P(x) \rightarrow Q(x))$ may be logically valid but evaluates to TRUE for all candidates (e.g., both ``Proved'' and ``Unknown'' in a deduction task).
The extraction process produces logically correct but epistemically vacuous constraints.

% ============================================================
% 4.5 C0 DIAGNOSTIC
% ============================================================
\subsection{$C_0$ as a Deployment Diagnostic}
\label{sec:c0}

We validate $C_0$ as a predictor of \sica{} effectiveness across all 18 experimental conditions, spanning $C_0$ from 2.6\% (QwQ-32B) to 70\% (Qwen-7B).

\paragraph{$C_0$ is necessary but not sufficient.}
All conditions with $C_0 < 5\%$ show $\Delta = 0.0$pp exactly, confirming that low diversity makes \sica{} structurally degenerate.
However, among conditions with $C_0 > 20\%$, the mean $|\Delta|$ is only 0.1pp with max $|\Delta| = 2.5$pp.
The correlation between $C_0$ and $\Delta$ across all 18 conditions is $r = 0.05$ ($p = 0.84$), confirming that diversity provides no gradient of improvement.

\paragraph{Practical interpretation.}
$C_0$ functions as a gate, not a dial.
Below 5\%, practitioners can skip symbolic aggregation entirely with confidence.
Above 5\%, $C_0$ provides no guidance on whether gains will materialize, because the binding constraint shifts to extraction quality, which $C_0$ does not measure.
Crucially, $C_0$ is computable from trace generation alone (Stage 1), before any constraint extraction, making it a zero-cost early-stopping criterion.

% ============================================================
% 4.6 TIEBREAK AND K-ABLATION
% ============================================================
\subsection{Tiebreak Analysis and K Ablation}
\label{sec:tiebreak}

\paragraph{Tiebreak is locally effective but globally irrelevant.}
On the 5 problems (out of 804) where SC produces an exact tie between candidates, \sica{}'s constraint-based tiebreak achieves 80\% accuracy vs.\ 50\% for random tiebreak and 58\% for alphabetical tiebreak.
However, ties occur in only 0.6\% of problems, contributing $<0.3$pp to overall accuracy, which is statistically negligible.

\paragraph{K ablation.}
We test $K \in \{6, 8, 11, 12, 16, 20\}$ on Qwen-14B ($n=60$).
\sica{} $\leq$ SC at all values of $K$, including odd $K=11$ where exact ties are structurally impossible.
The $K=11$ result ($\Delta = 0.0$pp) directly rules out the hypothesis that \sica{}'s value lies solely in tiebreaking.
Larger $K$ marginally degrades \sica{} (e.g., $K=20$: $\Delta = -1.7$pp), consistent with noise accumulation from additional non-discriminative constraints.

% ============================================================
% 4.7 ADDITIONAL BASELINES
% ============================================================
\subsection{Additional Baselines}
\label{sec:baselines}

\paragraph{Weighted SC.}
We compare four weighting schemes on the $n=804$ Qwen-14B data (Table~\ref{tab:weighted_sc}).
No scheme achieves significant improvement over uniform SC.
Inverse-length weighting shows the largest (but non-significant) edge at $+0.5$pp ($p=0.453$).
This eliminates the concern that simpler weighting could succeed where MAX-SAT fails: the bottleneck is signal quality, not aggregation sophistication.

\begin{table}[t]
\centering
\caption{Weighted SC variants ($n=804$, Qwen-14B). No weighting achieves significant improvement.}
\label{tab:weighted_sc}
\small
\begin{tabular}{lrr}
\toprule
\textbf{Weighting} & \textbf{Accuracy (\%)} & \textbf{$\Delta$ vs.\ SC (pp)} \\
\midrule
Uniform (SC) & 76.24 & -- \\
Inverse-length & 76.74 & +0.50 \\
Trace-length & 75.25 & $-$0.99 \\
Constraint-count & 75.25 & $-$0.99 \\
\bottomrule
\end{tabular}
\end{table}

\paragraph{Logic-LM $\times$ K baseline.}
On FOLIO ($n=204$), applying Logic-LM-style single-trace FOL verification $K$ times and voting produces accuracy identical to both SC (79.4\%) and \sica{} (79.4\%).
Whether using MAX-SAT (\sica{}), majority vote (SC), or single-trace verification (Logic-LM $\times K$), same-model extraction produces equivalent results.
The bottleneck is extraction quality, not aggregation mechanism.

% ============================================================
% 4.8 OVERHEAD
% ============================================================
\subsection{Overhead Analysis}
\label{sec:overhead}

\sica{} nearly doubles inference cost for negligible accuracy gain (Table~\ref{tab:overhead}).

\begin{table}[t]
\centering
\caption{Wall-time breakdown per problem ($n=804$, Qwen-14B). \sica{} adds $\sim$49\% overhead, dominated by LLM-based constraint extraction.}
\label{tab:overhead}
\small
\begin{tabular}{lrr}
\toprule
\textbf{Stage} & \textbf{Mean time (s)} & \textbf{Fraction (\%)} \\
\midrule
Trace generation & 31.6 & 51.3 \\
Constraint extraction & 29.0 & 47.0 \\
MAX-SAT solving & 0.003 & $<$0.01 \\
\midrule
Total & 61.6 & 100.0 \\
\bottomrule
\end{tabular}
\end{table}

The overhead is dominated by LLM calls for constraint extraction (47\% of total time), while MAX-SAT solving itself is negligible ($<$0.01\%).
The cost multiplier is $1.95\times$ compared to SC alone.
Combined with zero effectiveness, this definitively answers the practitioner question: \sica{} is not worth the overhead in any tested condition.

% ============================================================
% 5. DISCUSSION
% ============================================================
\section{Discussion}
\label{sec:discussion}

\paragraph{Anatomy of the capability trap.}
The capability trap has a precise mechanistic explanation.
Self-extracted constraints are logically correct in 75\% of cases and correctly evaluate against their source answer 52\% of the time.
However, only 12\% of constraints are \emph{discriminative}: they produce different satisfaction outcomes for different candidate answers.
The remaining 88\% are tautological or overly general, evaluating identically across all candidates.
When most constraints are non-discriminative, the MAX-SAT objective $\sum_c w_c \cdot \text{sat}(a, c)$ becomes proportional to the number of supporting traces, reducing the solver to a weighted vote.
The three-way equivalence of MAX-SAT, count-only, and random-weight aggregation (Section~\ref{sec:ablation}) is the direct consequence: constraint weights carry no information when the constraints themselves do not discriminate.

\paragraph{Why does self-extraction fail?}
The extraction process asks the LLM to formalize its own reasoning steps as FOL constraints.
The resulting formulas tend to capture general reasoning patterns (e.g., modus ponens applications, domain axioms) that are valid regardless of the specific answer.
This is structurally different from constraints designed to test specific answer candidates, which would require the model to reason \emph{about} its own uncertainty, a metacognitive capability that current LLMs lack in the self-extraction setting.

\paragraph{Within-family scaling.}
The Qwen scaling results (7B: $C_0=70\%$, 14B: $C_0=38\%$, 72B: $C_0=22\%$) reveal that capability and diversity trade off within a model family.
Weaker models produce more diverse (and thus more frequently disagreeing) outputs but extract worse constraints (precision 0.49 for 7B vs.\ 0.51 for 72B).
Stronger models extract marginally better constraints but produce insufficient diversity to exercise them.
This tradeoff is structural: it arises from the same capability that determines both answer consistency and constraint quality.

\paragraph{Relationship to published systems.}
Systems like DyLoS, Aristotle, FoVer, and LogicAgent achieve gains on benchmarks overlapping with ours.
The key distinction is architectural: these systems employ trained verifiers, external oracles, logic-complete decomposition, or multi-paradigm solver selection.
Our finding characterizes the specific failure mode of \emph{same-model self-extraction} pipelines, not symbolic reasoning broadly.
The capability trap may be breakable by decoupling the extraction and reasoning models, as these systems implicitly do.

\paragraph{Reasoning models bypass the need entirely.}
QwQ-32B, as a full reasoning model, achieves $C_0=2.6\%$, making \sica{} structurally degenerate.
This is consistent with the test-time compute scaling paradigm \citep{snell2025scaling,deepseekr1}: models that invest compute in extended chains achieve consistency through deliberation rather than aggregation.
For reasoning models, the question of symbolic aggregation is moot.

\paragraph{Practical recommendations.}
Practitioners considering symbolic aggregation over self-extracted constraints should:
(1) Compute $C_0$ from a small pilot ($n \geq 30$, $K=12$). If $C_0 < 5\%$, the method will degenerate to SC.
(2) Even if $C_0 > 20\%$, do not expect gains without evidence that the extraction model produces discriminative constraints (precision $\gg 0.50$).
(3) Consider the $\sim$2$\times$ overhead: SC achieves the same accuracy at half the cost.

% ============================================================
% 6. CONCLUSION
% ============================================================
\section{Conclusion}
\label{sec:conclusion}

We presented a systematic evaluation of symbolic constraint aggregation via self-extracted FOL constraints (SICA) across 3 reasoning domains, 5 model families, and 18 experimental conditions.
Our central finding is a \emph{capability trap}: self-extracted constraints achieve chance-level discriminative precision, causing MAX-SAT to degenerate to counting and providing no improvement over simple majority voting.
We validated this through constraint ablation (MAX-SAT $\equiv$ count-only $\equiv$ random-weight), per-problem precision analysis ($n=307$, mean precision $= 0.50$), and cross-condition C0 analysis ($r = 0.05$ between diversity and gain).
The multi-candidate rate $C_0$ provides a zero-cost diagnostic for practitioners: it reliably predicts when symbolic aggregation cannot help, potentially saving significant inference overhead ($\sim$2$\times$ cost).

\paragraph{Limitations.}
Our study is scoped to same-model self-extraction pipelines. Systems using trained verifiers, external oracles, or cross-model extraction may not share this limitation.
We test one MAX-SAT formulation; alternative symbolic aggregation strategies (e.g., SAT, SMT, answer set programming) could yield different results.
GSM8K and ARC-Challenge evaluations use smaller samples ($n \leq 200$); powered replications on these domains would strengthen conclusions.

% ============================================================
% REFERENCES
% ============================================================
\bibliographystyle{plainnat}

\begin{thebibliography}{30}

\bibitem[Wang et~al.(2023)]{wang2023selfconsistency}
Wang, X., Wei, J., Schuurmans, D., Le, Q., Chi, E., Narang, S., Chowdhery, A., and Zhou, D.
\newblock Self-Consistency Improves Chain of Thought Reasoning in Language Models.
\newblock In \emph{ICLR}, 2023.

\bibitem[Chen et~al.(2023)]{chen2023usc}
Chen, X., Aksitov, R., Alon, U., Ren, J., Xiao, K., Yin, P., Prakash, S., Sutton, C., Wang, X., and Zhou, D.
\newblock Universal Self-Consistency for Large Language Model Generation.
\newblock \emph{arXiv preprint arXiv:2311.17311}, 2023.

\bibitem[Pan et~al.(2023)]{pan2023logiclm}
Pan, L., Albalak, A., Wang, X., and Wang, W.Y.
\newblock Logic-LM: Empowering Large Language Models with Symbolic Solvers for Faithful Logical Reasoning.
\newblock In \emph{Findings of EMNLP}, 2023.

\bibitem[Olausson et~al.(2023)]{olausson2023linc}
Olausson, T.X., Gu, A., Lipkin, B., Zhang, C.E., Solar-Lezama, A., Tenenbaum, J.B., and Levy, R.
\newblock LINC: A Neurosymbolic Approach for Logical Reasoning by Combining Language Models with First-Order Logic Provers.
\newblock In \emph{EMNLP}, 2023.

\bibitem[Loo et~al.(2025)]{loo2025sc}
Loo, N., Hasanbeig, H., and Kamgarpour, M.
\newblock On the Scaling of Self-Consistency.
\newblock \emph{arXiv preprint}, 2025.

\bibitem[Taubenfeld et~al.(2025)]{taubenfeld2025weighted}
Taubenfeld, A., Shayegani, E., and Goldstein, T.
\newblock Confidence-Weighted Self-Consistency.
\newblock In \emph{ACL Findings}, 2025.

\bibitem[Li et~al.(2023)]{li2023diverse}
Li, Y., Lin, Z., Zhang, S., Fu, Q., Chen, B., Lou, J.-G., and Chen, W.
\newblock Making Language Models Better Reasoners with Step-Aware Verifier.
\newblock In \emph{ACL}, 2023.

\bibitem[DyLoS(2026)]{dylos2026}
DyLoS: Dynamic Logic-Solver Selection for Neuro-Symbolic Reasoning.
\newblock In \emph{EACL}, 2026.

\bibitem[Aristotle(2025)]{aristotle2025}
Aristotle: Logic-Complete Decomposition for Symbolic Reasoning.
\newblock In \emph{ACL}, 2025.

\bibitem[FoVer(2025)]{fover2025}
FoVer: Formal Verification with Trained Verifiers for Logical Reasoning.
\newblock In \emph{TACL}, 2025.

\bibitem[Cobbe et~al.(2021)]{cobbe2021gsm8k}
Cobbe, K., Kosaraju, V., Bavarian, M., Chen, M., Jun, H., Kaiser, L., Plappert, M., Tworek, J., Hilton, J., Nakano, R., Hesse, C., and Schulman, J.
\newblock Training Verifiers to Solve Math Word Problems.
\newblock \emph{arXiv preprint arXiv:2110.14168}, 2021.

\bibitem[Lightman et~al.(2023)]{lightman2023prm}
Lightman, H., Kosaraju, V., Burda, Y., Edwards, H., Baker, B., Lee, T., Leike, J., Schulman, J., Sutskever, I., and Alignment, K.
\newblock Let's Verify Step by Step.
\newblock In \emph{ICLR}, 2024.

\bibitem[Wang et~al.(2024)]{wang2024mathshepherd}
Wang, P., Li, L., Shao, Z., Xu, R., Dai, D., Li, Y., Chen, D., Wu, Y., and Sui, Z.
\newblock Math-Shepherd: Verify and Reinforce LLMs Step-by-step without Human Annotations.
\newblock In \emph{ACL}, 2024.

\bibitem[Snell et~al.(2025)]{snell2025scaling}
Snell, C., Lee, J., Xu, K., and Kumar, A.
\newblock Scaling LLM Test-Time Compute Optimally Can Be More Effective Than Scaling Model Parameters.
\newblock In \emph{ICLR (Oral)}, 2025.

\bibitem[DeepSeek-AI(2025)]{deepseekr1}
DeepSeek-AI.
\newblock DeepSeek-R1: Incentivizing Reasoning Capability in LLMs via Reinforcement Learning.
\newblock \emph{Nature}, 2025.

\bibitem[OpenAI(2024)]{openai2024o1}
OpenAI.
\newblock Learning to Reason with LLMs.
\newblock OpenAI Blog, 2024.

\bibitem[SATBench(2025)]{satbench2025}
SATBench: Benchmarking LLMs on Satisfiability Problems.
\newblock In \emph{EMNLP}, 2025.

\bibitem[LLM-CP(2025)]{llmcp2025}
LLM Meets Constraint Propagation.
\newblock In \emph{IJCAI}, 2025.

\bibitem[Liang et~al.(2023)]{liang2023debate}
Liang, T., He, Z., Jiao, W., Wang, X., Wang, Y., Wang, R., Yang, Y., Tu, Z., and Shi, S.
\newblock Encouraging Divergent Thinking in Large Language Models through Multi-Agent Debate.
\newblock \emph{arXiv preprint arXiv:2305.19118}, 2023.

\bibitem[Du et~al.(2023)]{du2023debate}
Du, Y., Li, S., Torralba, A., Tenenbaum, J.B., and Mordatch, I.
\newblock Improving Factuality and Reasoning in Language Models through Multiagent Debate.
\newblock In \emph{ICML}, 2024.

\bibitem[QwQ-Team(2024)]{qwq2024}
QwQ Team.
\newblock QwQ: Reflect Deeply on the Boundaries of the Unknown.
\newblock Qwen Blog, 2024.

\bibitem[Han et~al.(2022)]{han2022folio}
Han, S., Schoelkopf, H., Zhao, Y., Qi, Z., Riddell, M., Benson, L., Sun, L., Zubova, E., Qiao, Y., Burtell, M., Peng, D., Fan, J., Zhuo, T.Y., Jia, Y., and Radev, D.
\newblock FOLIO: Natural Language Reasoning with First-Order Logic.
\newblock \emph{arXiv preprint arXiv:2209.00840}, 2022.

\bibitem[Geva et~al.(2021)]{geva2021strategyqa}
Geva, M., Khashabi, D., Segal, E., Khot, T., Roth, D., and Berant, J.
\newblock Did Aristotle Use a Laptop? A Question Answering Benchmark with Implicit Reasoning Strategies.
\newblock In \emph{TACL}, 2021.

\bibitem[Clark et~al.(2018)]{clark2018arc}
Clark, P., Cowhey, I., Etzioni, O., Khot, T., Sabharwal, A., Schoenick, C., and Tafjord, O.
\newblock Think You Have Solved Question Answering? Try ARC, the AI2 Reasoning Challenge.
\newblock \emph{arXiv preprint arXiv:1803.05457}, 2018.

\bibitem[Cobbe et~al.(2021b)]{cobbe2021gsm8k2}
Cobbe, K., Kosaraju, V., Bavarian, M., Chen, M., Jun, H., Kaiser, L., Plappert, M., Tworek, J., Hilton, J., Nakano, R., Hesse, C., and Schulman, J.
\newblock Training Verifiers to Solve Math Word Problems.
\newblock \emph{arXiv preprint arXiv:2110.14168}, 2021.

\bibitem[Tafjord et~al.(2021)]{tafjord2021proofwriter}
Tafjord, O., Dalvi, B., and Clark, P.
\newblock ProofWriter: Generating Implications, Proofs, and Abductive Statements over Natural Language.
\newblock In \emph{Findings of ACL}, 2021.

\bibitem[Wei et~al.(2022)]{wei2022cot}
Wei, J., Wang, X., Schuurmans, D., Bosma, M., Xia, F., Chi, E., Le, Q., and Zhou, D.
\newblock Chain-of-Thought Prompting Elicits Reasoning in Large Language Models.
\newblock In \emph{NeurIPS}, 2022.

\bibitem[Qwen-Team(2024)]{qwen2024}
Qwen Team.
\newblock Qwen2.5: A Party of Foundation Models.
\newblock Technical Report, 2024.

\bibitem[Touvron et~al.(2023)]{touvron2023llama}
Touvron, H., Martin, L., Stone, K., et~al.
\newblock Llama 2: Open Foundation and Fine-Tuned Chat Models.
\newblock \emph{arXiv preprint arXiv:2307.09288}, 2023.

\bibitem[DeepSeek-AI(2024)]{deepseek2024}
DeepSeek-AI.
\newblock DeepSeek-V2: A Strong, Economical, and Efficient Mixture-of-Experts Language Model.
\newblock \emph{arXiv preprint arXiv:2405.04434}, 2024.

\end{thebibliography}

% ============================================================
% APPENDIX
% ============================================================
\appendix
\section{K Ablation Details}
\label{app:k_ablation}

Table~\ref{tab:k_ablation} shows \sica{} vs.\ SC across $K$ values on Qwen-14B ($n=60$, D012 subset).

\begin{table}[h]
\centering
\caption{K ablation (Qwen-14B, $n=60$). \sica{} $\leq$ SC at all K values.}
\label{tab:k_ablation}
\small
\begin{tabular}{rrrr}
\toprule
\textbf{K} & \textbf{\sica{} Acc (\%)} & \textbf{SC Acc (\%)} & \textbf{$\Delta$pp} \\
\midrule
6 & 73.3 & 73.3 & 0.0 \\
8 & 73.3 & 73.3 & 0.0 \\
11 & 73.3 & 73.3 & 0.0 \\
12 & 73.3 & 75.0 & $-$1.7 \\
16 & 71.7 & 73.3 & $-$1.7 \\
20 & 70.0 & 71.7 & $-$1.7 \\
\bottomrule
\end{tabular}
\end{table}

\section{Tiebreak Baseline Details}
\label{app:tiebreak}

On the $n=804$ Qwen-14B evaluation, only 5 problems (0.6\%) produce exact SC ties.
Table~\ref{tab:tiebreak} compares tiebreak strategies on these tied problems.

\begin{table}[h]
\centering
\caption{Tiebreak accuracy on 5 tied problems (Qwen-14B, $n=804$).}
\label{tab:tiebreak}
\small
\begin{tabular}{lr}
\toprule
\textbf{Tiebreak Strategy} & \textbf{Accuracy (\%)} \\
\midrule
\sica{} (constraint-based) & 80.0 \\
Alphabetical & 58.0 \\
Random & 50.0 \\
First-seen & 42.0 \\
\bottomrule
\end{tabular}
\end{table}

\sica{}'s constraint-based tiebreak is locally effective but contributes $<0.3$pp to overall accuracy due to the rarity of ties.

\section{Overhead Breakdown}
\label{app:overhead}

The $\sim$2$\times$ cost multiplier ($1.95\times$) of \sica{} over SC is dominated by constraint extraction LLM calls (47.0\% of total wall time).
MAX-SAT solving accounts for $<$0.01\% of computation.
This cost structure means that any optimization of \sica{}'s overhead must target the extraction step, not the solver.

\section{Within-Family Scaling}
\label{app:scaling}

\begin{table}[h]
\centering
\caption{Within-family scaling (Qwen, $n=60$, $K=12$). $C_0$ and constraint quality trade off, maintaining zero net gain.}
\label{tab:scaling}
\small
\begin{tabular}{lrrrrrr}
\toprule
\textbf{Model} & \textbf{SC (\%)} & \textbf{\sica{} (\%)} & \textbf{$\Delta$pp} & \textbf{$C_0$ (\%)} & \textbf{Extr.} & \textbf{Prec.} \\
\midrule
Qwen-7B & 58.3 & 58.3 & 0.0 & 70 & 0.78 & 0.49 \\
Qwen-14B & 73.3 & 73.3 & 0.0 & 38 & -- & -- \\
Qwen-72B & 80.0 & 80.0 & 0.0 & 22 & 0.94 & 0.51 \\
\bottomrule
\end{tabular}
\end{table}

The scaling data reveals a structural tradeoff: weaker models produce higher diversity ($C_0$) but lower constraint quality (extraction rate and precision), while stronger models improve constraint quality at the expense of diversity.
At no point along this spectrum does \sica{} outperform SC.

\section{Logic-LM Comparison}
\label{app:logiclm}

On FOLIO ($n=204$), Logic-LM-style single-trace FOL verification applied $K=12$ times with majority voting achieves 79.4\% accuracy, identical to both SC (79.4\%) and \sica{} (79.4\%).
Logic-LM single-trace accuracy is 78.9\%.
The equivalence across all three aggregation paradigms confirms that same-model extraction quality is the binding bottleneck.

\section{QwQ-32B Details}
\label{app:qwq}

QwQ-32B produces near-deterministic outputs across all domains: effective $C_0 = 2.6\%$ (only 1 out of 46 problems produces multiple candidates).
Per-domain results: AIME $\Delta=0.0$pp, ProofWriter $\Delta=0.0$pp, FOLIO $\Delta=0.0$pp.
\sica{} structurally degenerates to SC for full reasoning models.

\section{DeepSeek-R1-Distill-7B Details}
\label{app:deepseek}

Despite being distilled from a reasoning model, DeepSeek-R1-Distill-Qwen-7B ($n=60$) behaves as an instruction-scale model: $C_0=58.3\%$, comparable to other 7--8B instruction models.
The 7B capacity is insufficient to maintain deterministic reasoning chains.
$\Delta=-3.33$pp (n.s., McNemar $p=0.48$, 0 \sica{} wins vs.\ 2 SC wins among discordant pairs).
72 traces were truncated due to context length, potentially contributing to the negative direction.

\end{document}
